% ═══════════════════════════════════════════════════════════════
% ML-02a: Ohmsches Gesetz – MUSTERLÖSUNG
% Physik Klasse 8 | Doppelstunde 02a
% Struktur identisch mit AB-02a
% ═══════════════════════════════════════════════════════════════

\documentclass[11pt, a4paper]{article}

\usepackage[utf8]{inputenc}
\usepackage[T1]{fontenc}
\usepackage[ngerman]{babel}

\usepackage{helvet}
\renewcommand{\familydefault}{\sfdefault}

\usepackage{microtype}
\usepackage[margin=1.5cm, top=1.8cm, bottom=1.5cm]{geometry}
\usepackage{enumitem}
\usepackage{tabularx}
\usepackage{booktabs}
\usepackage{array}
\usepackage{multirow}
\usepackage{tikz}
\usetikzlibrary{calc}

\usepackage{amsmath, amssymb}
\usepackage{siunitx}
\sisetup{locale = DE, per-mode = symbol}

\usepackage{fancyhdr}
\usepackage{lastpage}
\usepackage{needspace}

\usepackage{xcolor}
\usepackage{tcolorbox}
\tcbuselibrary{skins, breakable}

% Farben
\definecolor{physik}{HTML}{2c5aa0}
\definecolor{hellgrau}{HTML}{F5F5F5}
\definecolor{mittelgrau}{HTML}{888888}
\definecolor{loesung}{HTML}{c62828}

% Variablen
\newcommand{\fach}{Physik}
\newcommand{\thema}{Das Ohmsche Gesetz}
\newcommand{\klasse}{8}
\newcommand{\maxpunkte}{47}

\colorlet{fachfarbe}{physik}

% Kopf- und Fußzeile
\pagestyle{fancy}
\fancyhf{}
\renewcommand{\headrulewidth}{0.4pt}
\lhead{\small \fach{} Klasse \klasse{} | \thema{} -- \textcolor{loesung}{\textbf{MUSTERLÖSUNG}}}
\rhead{\small}
\cfoot{\small\textcolor{mittelgrau}{Seite \thepage{} von \pageref{LastPage}}}

% Boxen
\newtcolorbox{merkkasten}[1][]{
    colback=hellgrau,
    colframe=hellgrau,
    boxrule=0pt,
    arc=0pt,
    left=8pt, right=8pt, top=8pt, bottom=6pt,
    borderline north={2pt}{0pt}{fachfarbe},
    before skip=8pt, after skip=8pt,
    #1
}

\newtcolorbox{formelkasten}{
    colback=white,
    colframe=fachfarbe,
    boxrule=1.5pt,
    arc=0pt,
    left=8pt, right=8pt, top=6pt, bottom=6pt,
    before skip=6pt, after skip=6pt
}

\newtcolorbox{hinweiskasten}{
    colback=white,
    colframe=white,
    boxrule=0pt,
    arc=0pt,
    left=8pt, right=4pt, top=4pt, bottom=4pt,
    borderline west={3pt}{0pt}{fachfarbe}
}

% Befehle
\newcommand{\luecke}[1]{\underline{\hspace{#1}}}
\newcommand{\loes}[1]{\textcolor{loesung}{\textbf{#1}}}
\newcommand{\punktebox}[1]{\hfill\small\textcolor{mittelgrau}{[\hspace{0.8em}/#1]}}

\newcommand{\aufgabe}[2]{%
    \needspace{4\baselineskip}%
    \vspace{8pt}%
    \noindent\textbf{\textcolor{fachfarbe}{Aufgabe #1}} #2%
    \vspace{4pt}%
    \nopagebreak[4]%
}

\newcommand{\operator}[1]{\textbf{#1}}

\newlist{teil}{enumerate}{2}
\setlist[teil,1]{label=\alph*), leftmargin=1.8em, itemsep=4pt, parsep=0pt, topsep=4pt}

\setlength{\parindent}{0pt}
\setlength{\parskip}{3pt}

\begin{document}

% ═══════════════════════════════════════════════════════════════
% HEADER
% ═══════════════════════════════════════════════════════════════

\begin{center}
    {\Large\bfseries\textcolor{fachfarbe}{\thema}}\\[0.2em]
    {\small Arbeitsblatt | \fach{} Klasse \klasse{} | \textcolor{loesung}{\textbf{MUSTERLÖSUNG}}}
\end{center}
\vspace{-0.5em}
\hrule
\vspace{0.8em}

% ═══════════════════════════════════════════════════════════════
% KASTEN 1: WASSERMODELL
% ═══════════════════════════════════════════════════════════════

\begin{merkkasten}
\textbf{Kasten 1: Die Wasser-Analogie}

\vspace{0.3em}
Elektrischer Strom verhält sich ähnlich wie fließendes Wasser:

\vspace{0.5em}
\begin{center}
\begin{tabular}{|>{\raggedright\arraybackslash}p{4.5cm}|p{5cm}|}
\hline
\textbf{Wasserkreislauf} & \textbf{Stromkreis} \\
\hline
Wasserdruck (durch Pumpe) & \loes{Spannung $U$} \\[0.6em]
\hline
Wassermenge pro Sekunde & \loes{Stromstärke $I$} \\[0.6em]
\hline
Engstelle im Rohr & \loes{Widerstand $R$} \\[0.6em]
\hline
\end{tabular}
\end{center}
\end{merkkasten}

% ═══════════════════════════════════════════════════════════════
% AUFGABE 1: MESSWERTE
% ═══════════════════════════════════════════════════════════════

\aufgabe{1}{$(*)$ – Messwerte aufnehmen} \punktebox{6}

\operator{Nimm} mithilfe der Simulation Messwerte auf. \operator{Trage} Spannung und Stromstärke in die Tabelle ein.

\vspace{0.5em}
\begin{center}
\begin{tabular}{|c|c||c|}
\hline
\textbf{Spannung $U$ in V} & \textbf{Stromstärke $I$ in A} & \textbf{Widerstand $R$ in $\Omega$} \\
& & \textit{\footnotesize (später ergänzen!)} \\
\hline
2 & \loes{0,20} & \loes{10} \\[0.5em]
\hline
4 & \loes{0,40} & \loes{10} \\[0.5em]
\hline
6 & \loes{0,60} & \loes{10} \\[0.5em]
\hline
8 & \loes{0,80} & \loes{10} \\[0.5em]
\hline
10 & \loes{1,00} & \loes{10} \\[0.5em]
\hline
12 & \loes{1,20} & \loes{10} \\[0.5em]
\hline
\end{tabular}
\end{center}

% ═══════════════════════════════════════════════════════════════
% AUFGABE 2: KENNLINIE ZEICHNEN
% ═══════════════════════════════════════════════════════════════

\aufgabe{2}{$(*)$ – Kennlinie zeichnen} \punktebox{8}

\begin{teil}
    \item \operator{Trage} deine Messpunkte in das Koordinatensystem ein.
    \item \operator{Zeichne} eine \textbf{Ausgleichskurve} durch die Punkte (mit Lineal!).
    \item \operator{Verlängere} die Kurve bis zum Rand des Koordinatensystems.
\end{teil}

\vspace{0.3em}
\begin{center}
\begin{tikzpicture}[scale=0.8]
    % Achsen
    \draw[thick, ->] (0,0) -- (14,0) node[right] {$U$ in V};
    \draw[thick, ->] (0,0) -- (0,8.5) node[above] {$I$ in A};

    % Gitternetz
    \draw[gray!30, very thin] (0,0) grid[step=1] (13,7);

    % Beschriftung x-Achse (0 bis 12 V)
    \foreach \x in {0,2,4,6,8,10,12}
        \draw (\x,0) -- (\x,-0.15) node[below] {\small\x};

    % Beschriftung y-Achse (0 bis 1,4 A, Schritt 0,2)
    % 1 Grid-Einheit = 0,2 A, also y=5 entspricht 1,0 A
    \foreach \y/\ylabel in {0/0, 1/0{,}2, 2/0{,}4, 3/0{,}6, 4/0{,}8, 5/1{,}0, 6/1{,}2, 7/1{,}4}
        \draw (0,\y) -- (-0.15,\y) node[left] {\small\ylabel};

    % LÖSUNG: Messpunkte ohmscher Widerstand (blau)
    % I = U/10, also bei 0,2 A/Grid: y = U/10 * 5 = U/2
    \foreach \u in {2, 4, 6, 8, 10, 12}
        \fill[fachfarbe] (\u, \u/2) circle (3pt);

    % LÖSUNG: Ausgleichsgerade ohmscher Widerstand (blau)
    \draw[fachfarbe, thick] (0,0) -- (13,6.5);

    % LÖSUNG: Messpunkte Glühlampe (rot/orange)
    % I-Werte aus Simulation: 0.32, 0.50, 0.62, 0.72, 0.80, 0.86
    % Bei 0,2 A/Grid: y = I * 5
    \fill[loesung] (2, 1.6) circle (3pt);   % 0.32 * 5 = 1.6
    \fill[loesung] (4, 2.5) circle (3pt);   % 0.50 * 5 = 2.5
    \fill[loesung] (6, 3.1) circle (3pt);   % 0.62 * 5 = 3.1
    \fill[loesung] (8, 3.6) circle (3pt);   % 0.72 * 5 = 3.6
    \fill[loesung] (10, 4.0) circle (3pt);  % 0.80 * 5 = 4.0
    \fill[loesung] (12, 4.3) circle (3pt);  % 0.86 * 5 = 4.3

    % LÖSUNG: Kennlinie Glühlampe (gekrümmt, rot)
    \draw[loesung, thick] plot[smooth] coordinates {(0,0) (2,1.6) (4,2.5) (6,3.1) (8,3.6) (10,4.0) (12,4.3) (13,4.5)};

    % Legende
    \node[fachfarbe, right] at (10, 6.5) {\small Ohmscher Widerstand};
    \node[loesung, right] at (10, 5.5) {\small Glühlampe};
\end{tikzpicture}
\end{center}

\vspace{0.3em}
\textit{Was fällt dir auf?} Die Messpunkte liegen auf einer \loes{Geraden} durch den \loes{Ursprung}.

% ═══════════════════════════════════════════════════════════════
% SEITE 2
% ═══════════════════════════════════════════════════════════════

\newpage

% ═══════════════════════════════════════════════════════════════
% KASTEN 2: OHMSCHES GESETZ
% ═══════════════════════════════════════════════════════════════

\begin{merkkasten}
\textbf{Kasten 2: Das Ohmsche Gesetz}

\vspace{0.4em}
\textbf{Deine Entdeckung:} Die Kennlinie ist eine Gerade durch den Ursprung!

\vspace{0.3em}
Das bedeutet: Die \loes{Stromstärke $I$} ist \textbf{proportional} zur \loes{Spannung $U$}.

\vspace{0.3em}
In Formelzeichen: \hspace{1em} \fbox{\large $I \sim U$} \hspace{1em} (bei konstantem Widerstand)

\vspace{0.6em}
\hrule
\vspace{0.6em}

\textbf{Das Ohmsche Gesetz} gibt dieser Proportionalitätskonstante einen Namen:

\begin{formelkasten}
\begin{center}
\large
$R = \dfrac{U}{I}$ \hspace{2cm} \textbf{Widerstand} $=$ \textbf{Spannung} geteilt durch \textbf{Stromstärke}
\end{center}
\end{formelkasten}

\vspace{0.3em}
\textbf{Einheit:} $[R] = 1\,\Omega$ (Ohm) $= 1\,\dfrac{\text{V}}{\text{A}}$ (Volt pro Ampere)

\vspace{0.6em}
\textbf{Die drei Formen} (durch Umstellen):

\vspace{0.3em}
\begin{center}
\begin{tabular}{|c|c|c|}
\hline
\textbf{Widerstand} & \textbf{Spannung} & \textbf{Stromstärke} \\[0.3em]
\hline
& & \\[0.3em]
$R = \dfrac{U}{I}$ & $U = $ \loes{$R \cdot I$} & $I = $ \loes{$U / R$} \\[0.8em]
\hline
\end{tabular}
\end{center}

\vspace{0.5em}
\begin{minipage}[t]{0.5\linewidth}
\textbf{Formeldreieck} (Merkhilfe):

\begin{center}
\begin{tikzpicture}[scale=0.6]
    \draw[thick, fachfarbe] (0,0) -- (4,0) -- (2,3) -- cycle;
    \draw[fachfarbe] (0.5,1.5) -- (3.5,1.5);
    \node at (2,2.2) {\large $U$};
    \node at (0.9,0.6) {\large $R$};
    \node at (3.1,0.6) {\large $I$};
    \node at (2,0.6) {\large $\cdot$};
\end{tikzpicture}
\end{center}
\end{minipage}%
\begin{minipage}[t]{0.5\linewidth}
\small
\textit{Anwendung:}\\[0.3em]
Was du suchst, deckst du ab!\\[0.2em]
• $U$ abdecken $\rightarrow$ $U = R \cdot I$\\[0.1em]
• $R$ abdecken $\rightarrow$ $R = U / I$\\[0.1em]
• $I$ abdecken $\rightarrow$ $I = U / R$
\end{minipage}
\end{merkkasten}

% ═══════════════════════════════════════════════════════════════
% AUFGABE 3: R BERECHNEN
% ═══════════════════════════════════════════════════════════════

\aufgabe{3}{$(*)$ – Widerstand berechnen} \punktebox{4}

\operator{Berechne} jetzt den Widerstand $R$ für alle deine Messwerte aus Aufgabe 1.

\vspace{0.3em}
\begin{hinweiskasten}
\textit{Gehe zurück zu Aufgabe 1 und fülle die rechte Spalte aus!}
\end{hinweiskasten}

\vspace{0.3em}
Was fällt dir auf? \loes{Der Widerstand $R$ ist bei allen Messungen gleich (konstant) = 10 $\Omega$}

% ═══════════════════════════════════════════════════════════════
% AUFGABE 4: BERECHNUNGEN ÜBEN
% ═══════════════════════════════════════════════════════════════

\aufgabe{4}{$(*)$/$(**)$ – Berechnungen mit dem Ohmschen Gesetz} \punktebox{12}

\begin{teil}
    \item $(*)$ Gegeben: $U = 12\,\text{V}$, $I = 3\,\text{A}$. \operator{Berechne} $R$. \hfill (2 BE)

    \vspace{0.3em}
    Formel: \loes{$R = U/I$} \hspace{0.5cm} Rechnung: \loes{$R = 12\,\text{V} / 3\,\text{A}$} \hspace{0.5cm} $R = $ \loes{4 $\Omega$}

    \vspace{0.6em}
    \item $(*)$ Gegeben: $R = 5\,\Omega$, $I = 2\,\text{A}$. \operator{Berechne} $U$. \hfill (2 BE)

    \vspace{0.3em}
    Formel: \loes{$U = R \cdot I$} \hspace{0.5cm} Rechnung: \loes{$U = 5\,\Omega \cdot 2\,\text{A}$} \hspace{0.5cm} $U = $ \loes{10 V}

    \vspace{0.6em}
    \item $(*)$ Gegeben: $U = 24\,\text{V}$, $R = 8\,\Omega$. \operator{Berechne} $I$. \hfill (2 BE)

    \vspace{0.3em}
    Formel: \loes{$I = U/R$} \hspace{0.5cm} Rechnung: \loes{$I = 24\,\text{V} / 8\,\Omega$} \hspace{0.5cm} $I = $ \loes{3 A}

    \vspace{0.6em}
    \item $(**)$ Ein Widerstand hat $R = 50\,\Omega$. Bei welcher Spannung fließen $I = 0{,}4\,\text{A}$? \hfill (3 BE)

    \textcolor{loesung}{$U = R \cdot I = 50\,\Omega \cdot 0{,}4\,\text{A} = \textbf{20 V}$}

    \vspace{0.5em}
    \item $(**)$ Gegeben: $R = 100\,\Omega$, $U = 20\,\text{V}$. Berechne $I$ \textbf{in mA}. \hfill (3 BE)

    \textcolor{loesung}{$I = U / R = 20\,\text{V} / 100\,\Omega = 0{,}2\,\text{A} = \textbf{200 mA}$}
\end{teil}

% ═══════════════════════════════════════════════════════════════
% SEITE 3
% ═══════════════════════════════════════════════════════════════

\newpage

% ═══════════════════════════════════════════════════════════════
% AUFGABE 5: GLÜHLAMPE
% ═══════════════════════════════════════════════════════════════

\aufgabe{5}{$(**)$ – Die Glühlampe: Ein besonderer Widerstand} \punktebox{11}

\begin{teil}
    \item \operator{Nimm} mit der zweiten Simulation Messwerte für eine \textbf{Glühlampe} auf. \operator{Berechne} $R = U/I$. \hfill (4 BE)

    \vspace{0.3em}
    \begin{center}
    \small
    \begin{tabular}{|c|c|c|c|c|c|c|}
    \hline
    $U$ in V & 2 & 4 & 6 & 8 & 10 & 12 \\
    \hline
    $I$ in A & \loes{0,32} & \loes{0,50} & \loes{0,62} & \loes{0,72} & \loes{0,80} & \loes{0,86} \\[0.5em]
    \hline
    $R$ in $\Omega$ & \loes{6,3} & \loes{8,0} & \loes{9,7} & \loes{11,1} & \loes{12,5} & \loes{14,0} \\[0.5em]
    \hline
    \end{tabular}
    \end{center}

    \vspace{0.3em}
    Was fällt dir bei den $R$-Werten auf? \loes{$R$ ist nicht konstant – er steigt mit der Spannung!}

    \vspace{0.5em}
    \item \operator{Zeichne} die Kennlinie der Glühlampe in \textbf{einer anderen Farbe} in das Diagramm auf Seite 1 ein. \hfill (2 BE)

    \textcolor{loesung}{(Siehe rote Kurve im Diagramm auf Seite 1)}

    \vspace{0.5em}
    \item \operator{Vergleiche} die beiden Kennlinien. Was ist anders? \hfill (2 BE)

    \textcolor{loesung}{Die Kennlinie der Glühlampe ist eine \textbf{Kurve} (keine Gerade). Sie verläuft \textbf{flacher} als die ohmsche Gerade. Der Strom steigt langsamer als erwartet.}

    \vspace{0.3em}
    \item \operator{Erkläre}, warum die Kennlinie der Glühlampe gekrümmt ist. \hfill (3 BE)

    \vspace{0.3em}
    \textit{Fülle die Lücken:}

    \vspace{0.4em}
    Bei höherer Spannung fließt mehr \loes{Strom}. Die Elektronen stoßen häufiger mit

    \vspace{0.3em}
    den \loes{Atomrümpfen}. Dadurch beginnen diese stärker zu \loes{schwingen}.

    \vspace{0.3em}
    Der Draht wird \loes{heiß}. Je heißer der Draht, desto \loes{größer} ist der

    \vspace{0.3em}
    Widerstand. Deshalb steigt der Strom \loes{weniger} stark als erwartet.

    \vspace{0.4em}
    \hfill\textit{\small (Wörter: Strom, Atomrümpfen, schwingen, heiß, größer, weniger)}
\end{teil}

% ═══════════════════════════════════════════════════════════════
% AUFGABE 6: TRANSFER
% ═══════════════════════════════════════════════════════════════

\aufgabe{6}{$(***)$ – Anwendungen} \punktebox{6}

\begin{teil}
    \item Ein Föhn ist mit $230\,\text{V}$ und $10\,\text{A}$ beschriftet. \operator{Berechne} $R$. \hfill (2 BE)

    \textcolor{loesung}{$R = U / I = 230\,\text{V} / 10\,\text{A} = \textbf{23 $\Omega$}$}

    \vspace{0.3em}
    \item Verlängerungskabel werden manchmal warm. \operator{Erkläre} warum. \hfill (2 BE)

    \textcolor{loesung}{Das Kabel hat selbst einen kleinen Widerstand. Bei hohem Strom (viele Geräte) stoßen die Elektronen häufiger mit den Atomrümpfen im Kupferdraht. Dabei wird Energie in Wärme umgewandelt $\rightarrow$ das Kabel wird warm.}

    \vspace{0.3em}
    \item Lisa berechnet $R = 8\,\Omega$. Richtig wäre $R = 4\,\Omega$. \operator{Erkläre} ihren Fehler. \hfill (2 BE)

    \textcolor{loesung}{Lisa hat wahrscheinlich $U$ und $I$ vertauscht, also $I/U$ statt $U/I$ gerechnet. Oder sie hat sich beim Teilen verrechnet (z.B. falsch dividiert).}
\end{teil}

% ═══════════════════════════════════════════════════════════════
% KASTEN 3: ZUSAMMENFASSUNG
% ═══════════════════════════════════════════════════════════════

\vspace{0.3em}
\begin{merkkasten}
\textbf{Kasten 3: Zusammenfassung}

\vspace{0.3em}
\begin{itemize}[leftmargin=1.5em, itemsep=2pt]
    \item Ohmsches Gesetz: \loes{$I$} $\sim$ \loes{$U$} (Proportionalität, bei konstantem $R$ und konstanter Temperatur)
    \item Formeln: $R =$ \loes{$U/I$}, \hspace{0.2cm} $U =$ \loes{$R \cdot I$}, \hspace{0.2cm} $I =$ \loes{$U/R$}
    \item Einheit: $[R] = 1\,\Omega = $ \loes{1 V/A}
    \item Ohmsche Widerstände: \loes{lineare} Kennlinie ($R$ konstant)
    \item Glühlampen: \loes{gekrümmte} Kennlinie ($R$ steigt mit \loes{Temperatur})
\end{itemize}
\end{merkkasten}

% ═══════════════════════════════════════════════════════════════
% PUNKTE
% ═══════════════════════════════════════════════════════════════

\vspace{0.8em}
\hrule
\vspace{0.5em}

\begin{center}
\small
\begin{tabular}{ccc}
\textbf{$\star$ Basis:} Aufg. 1, 2, 3, 4a-c & \textbf{$\star\star$ Standard:} + Aufg. 4d-e, 5 & \textbf{$\star\star\star$ Erweitert:} + Aufg. 6 \\
(24 BE erreichbar) & (41 BE erreichbar) & (47 BE erreichbar)
\end{tabular}
\end{center}

\vspace{0.3em}
\hfill\textbf{Punkte:} \luecke{1.5cm} / \maxpunkte

\end{document}
